\section*{Chapitre \ref{ch:m2:corporate-bonds} --- Obligations d'entreprise}

\begin{solution}{exo:m2:corporate-bonds:1}
\omterm{def:m2:corporate-bonds:spreads}{G-spread} $6.10 - 4.68 = 142$ points de base~; \omterm{def:m2:corporate-bonds:spreads}{I-spread} $6.10 - 4.40 = 170$.
\end{solution}

\begin{solution}{exo:m2:corporate-bonds:2}
Il a une notation scindée~: \omterm{def:m2:corporate-bonds:rating}{investment grade} auprès d'une agence, \omterm{def:m2:corporate-bonds:rating}{high yield} auprès de l'autre.
Son classement en \omterm{def:m2:corporate-bonds:rating}{investment grade} dépend de la règle de l'indice ou du mandat (certains retiennent
la plus basse notation, d'autres la médiane des trois). Selon la définition de la BCE, son émetteur est
un \omterm{def:m2:corporate-bonds:rating}{fallen angel}, déclassé en \omterm{def:m2:corporate-bonds:rating}{high yield} par au moins une grande agence.
\end{solution}

\begin{solution}{exo:m2:corporate-bonds:3}
$-6 \times 0.0080 \times 100 = -4.8$ pour 100.
\end{solution}

\begin{solution}{exo:m2:corporate-bonds:4}
Les deux mesurent l'écart entre le prix de l'obligation et sa valeur sur la courbe des swaps, l'un sous forme de spread courant au pair, l'autre
sous forme de décalage de chaque taux d'actualisation~; près du pair et pour des spreads modérés,
ils concordent au premier ordre, et les 2.2 points de base résultent de différences de capitalisation et de pondération.
Elles s'écartent à mesure que le spread croît et que le prix
s'éloigne du pair~: recalculée à 85, l'obligation a un \omterm{def:m2:corporate-bonds:spreads}{Z-spread} de 401.6 points de base et un
\omterm{def:m2:corporate-bonds:spreads}{spread d'asset swap} de 375.5.
\end{solution}

\begin{solution}{exo:m2:corporate-bonds:5}
$0.0140/0.60 = 2.33\%$ par an. Le reste du spread, quelles que soient les défaillances effectives
en deçà de ce niveau, rémunère l'incertitude sur le défaut et la perte,
l'illiquidité et, sur certains marchés, les impôts et le coût de détention de l'obligation.
\end{solution}

\begin{solution}{exo:m2:corporate-bonds:6}
Les investisseurs doivent absorber d'un coup une forte émission nouvelle et s'engager avant que le carnet soit
complet~; l'émetteur paie un peu pour le placer vite et en gros volume. Si la \omterm{def:m2:corporate-bonds:nic}{prime d'émission}
disparaît à mesure que les obligations remontent vers la courbe de l'émetteur, les souscripteurs
à l'émission en bénéficient, et l'émetteur l'a déjà versée.
\end{solution}

\begin{solution}{exo:m2:corporate-bonds:7}
À 105, le rendement est de 4.6548~\%~: \omterm{def:m2:corporate-bonds:spreads}{G-spread} 15.5 points de base, \omterm{def:m2:corporate-bonds:spreads}{I-spread} 40.5, \omterm{def:m2:corporate-bonds:spreads}{Z-spread} 41.2,
\omterm{def:m2:corporate-bonds:spreads}{spread d'asset swap} 43.2.
\end{solution}

\begin{solution}{exo:m2:corporate-bonds:8}
Une partie du spread de l'\omterm{def:m2:the-rates-options-market:swaption}{obligation remboursable par anticipation} rémunère l'option d'achat que son détenteur a vendue à l'émetteur~: son
\omterm{def:m2:corporate-bonds:oas}{spread ajusté des options} (OAS), à savoir la composante de crédit, est plus faible. Dans
les termes de l'\cref{ex:m2:corporate-bonds:four}, un \omterm{def:m2:corporate-bonds:spreads}{Z-spread} de 141 points de base est devenu un OAS
d'environ 61 une fois qu'une option d'achat à trois ans a été valorisée. Il faut comparer des OAS
entre elles, non des \omterm{def:m2:corporate-bonds:spreads}{Z-spread} entre eux.
\end{solution}

\begin{solution}{pb:m2:corporate-bonds:1}
\textbf{1.} $6 \times 1.50\% = 9.0$ pour 100.
\textbf{2.} $6 \times 0.80\% = 4.8$ pour 100, soit 53~\% de la baisse.
\textbf{3.} Les investisseurs lisent les mêmes résultats et suivent les mêmes évolutions de l'effet de levier et de la perspective
que les agences, et agissent en continu~; les agences agissent par comité et
par annonce.
\textbf{4.} Vendre tôt, ou réduire sa position, si son mandat ou sa politique de risque traite une
perspective négative à BBB$-$ comme un signal d'alerte.
\textbf{5.} $6 \times 0.70\% = 4.2$ pour 100.
\textbf{6.} 50 points de base au-dessus du spread juste de 350~: 3.0 pour 100.
\textbf{7.} Environ 1.5~million de dollars sur 50~millions de dollars, vendus à 3 points sous la fair value.
\textbf{8.} Les fonds \omterm{def:m2:corporate-bonds:rating}{high yield} et les investisseurs pouvant détenir du \omterm{def:m2:corporate-bonds:rating}{high yield}~; ils doivent absorber
d'un coup une offre importante et la financer en vendant d'autres obligations, et ils savent que les vendeurs
sont contraints.
\textbf{9.} Les règles des fournisseurs d'indices diffèrent, si bien que tous les indices n'excluent pas l'obligation le
même jour~; les fonds qui répliquent un échantillon de l'indice et peuvent garder une obligation
quelque temps étalent leurs ventes, ce qui allège la pression.
\textbf{10.} Le dépassement, si le spread revient à 350~: 3 points, soit 1.5~million de dollars,
plus le surcroît de revenu de coupon.
\textbf{11.} 3.0 pour 100, de 400 à 350.
\textbf{12.} De nouvelles dégradations ou un défaut, et un spread qui ne se redresse pas.
\textbf{13.} Que la plupart des \omterm{def:m2:corporate-bonds:rating}{fallen angels} survivent (près d'un quart est revenu en
\omterm{def:m2:corporate-bonds:rating}{investment grade}, près de la moitié est restée en \omterm{def:m2:corporate-bonds:rating}{high yield}), mais 12~\% ont fait défaut en vingt ans.
\textbf{14.} Environ $4.00\% \times 0.25 = 1.0$ pour 100 sur trois mois, au titre du
carry de spread, en plus du gain de prix.
\textbf{15.} L'acheteur est rémunéré pour prendre le risque que les vendeurs doivent se décharger, et le
spread peut encore s'élargir avant de se redresser.
\textbf{16.} Parce que les mandats d'investissement, les indices et les réglementations sont rédigés en
notations~: le seuil décide qui peut détenir l'obligation, quel que soit son prix.
\textbf{17.} En offrant à un acheteur les obligations récemment déclassées, dans certaines limites, il
réduirait le dépassement, comme le visait la disposition BB$-$/Ba3 du dispositif de 2020.
\textbf{18.} Conserver un \omterm{def:m2:corporate-bonds:rating}{fallen angel} pendant une période, ou jusqu'à une limite, plutôt que de vendre
d'un coup.
\textbf{19.} Résultat nommé~: \emph{la décote du \omterm{def:m2:corporate-bonds:rating}{fallen angel}}~: les ventes forcées ont poussé
les obligations 50 points de base au-delà de la fair value, soit 3 points pour 100, et 1.5~million de dollars sur la
position, qu'un acheteur a regagné en trois mois.
\textbf{20.} Parce que de nombreux détenteurs doivent vendre au même moment à des acheteurs qui ont besoin
d'une décote pour les absorber.
\end{solution}

\begin{solution}{iq:m2:corporate-bonds:1}
Le \omterm{def:m2:corporate-bonds:spreads}{G-spread} est le rendement excédant le rendement souverain interpolé~; le \omterm{def:m2:corporate-bonds:spreads}{Z-spread} est le
spread constant sur la courbe zéro qui réévalue tous les flux~; le \omterm{def:m2:corporate-bonds:spreads}{spread d'asset swap}
est ce qu'un investisseur gagne au-dessus du taux variable en échangeant au pair
les coupons de l'obligation. Elles diffèrent par le \omterm{def:m2:interest-rate-swaps:spread}{spread de swap}, la forme de la courbe et le
prix de l'obligation s'éloignant du pair.
\iqlookfor{définitions et raisons de leurs écarts.}
\end{solution}

\begin{solution}{iq:m2:corporate-bonds:2}
Un émetteur déclassé de l'\omterm{def:m2:corporate-bonds:rating}{investment grade} au \omterm{def:m2:corporate-bonds:rating}{high yield}. Ses obligations sortent
des indices \omterm{def:m2:corporate-bonds:rating}{investment grade} et de nombreux mandats, si bien que les détenteurs doivent vendre à un groupe plus restreint de
acheteurs qui exigent une décote~; le marché réévalue aussi le crédit avant le déclassement.
Les données montrent qu'une grande partie du mouvement précède l'événement et qu'un rétablissement partiel suit.
\iqlookfor{le mécanisme des ventes forcées et le calendrier.}
\end{solution}

\begin{solution}{iq:m2:corporate-bonds:3}
Estimer la perte attendue à partir des probabilités de défaut et des \omterm{def:m2:corporate-bonds:covenant}{taux de recouvrement}, historiques ou
issus d'un modèle, par notation et par échéance~; la retrancher du spread au-dessus de la
courbe sans risque~; le reliquat rémunère le risque et l'illiquidité, qui peuvent être
rattachés à la volatilité de marché, à des mesures de liquidité et à la capacité des investisseurs. Vérifier le résultat
pour toutes les notations et sur l'ensemble du cycle.
\iqlookfor{perte attendue face à la prime, et sa variation dans le temps.}
\end{solution}

\begin{solution}{iq:m2:corporate-bonds:4}
Partir des autres obligations de l'émetteur et de ses CDS, interpolés à l'échéance, et des
émetteurs comparables~; ajuster selon la taille, le coupon, les \omterm{def:m2:corporate-bonds:covenant}{covenants} et la liquidité~; vérifier
les transactions récentes et les cotations sur les plateformes~; puis valoriser le risque de détenir 50~millions de dollars jusqu'à
pouvoir la couvrir ou la vendre, ce qui élargit l'écart de prix selon la taille et l'illiquidité.
\iqlookfor{valeur relative selon la courbe et les comparables, plus un surcoût lié à la taille.}
\end{solution}

\begin{solution}{iq:m2:corporate-bonds:5}
Le \omterm{def:m2:corporate-bonds:spreads}{Z-spread} inclut la rémunération de l'option d'achat vendue par le détenteur~; l'OAS retire
la valeur de l'option, calculée à l'aide d'un modèle de taux, et conserve la composante de crédit et de liquidité.
Il faut comparer les OAS d'une obligation à l'autre lorsque les options diffèrent~; les \omterm{def:m2:corporate-bonds:spreads}{Z-spreads} ne se comparent qu'entre obligations bullet.
\iqlookfor{valeur de l'option incluse dans le spread.}
\end{solution}

\begin{solution}{iq:m2:corporate-bonds:6}
Maintenir à jour les courbes (Treasury, swap, taux zéro) à chaque événement de marché~; conserver
pour chaque obligation le calendrier des flux précalculé~; recalculer les prix et les spreads par lots vectorisés,
avec des solveurs de racines initialisés à partir de la dernière valeur~; paralléliser le calcul
entre obligations et ne recalculer que celles dont le prix ou la courbe a changé~; valider
les résultats par rapport à une implémentation de référence et signaler les cas de non-convergence.
\iqlookfor{précalcul, initialisations à chaud, mises à jour incrémentales.}
\end{solution}
